\chapter{Negación (\textit{niet} / \textit{geen})}
\section{Objetivos}
\begin{itemize}[leftmargin=*]
	\item Diferenciar el uso de \textit{niet} y \textit{geen} en nivel A1.
	\item Colocar la negación en la posición correcta en frases simples.
\end{itemize}

\section{Contenidos clave}
\begin{itemize}[leftmargin=*]
	\item \textit{geen}: negar sustantivos indefinidos o sin artículo.
	\item \textit{niet}: negar verbos, adjetivos, adverbios, preposiciones, elementos específicos.
	\item Orden básico: \textit{S + V2 + \ldots{} + niet} y combinaciones con partículas separables (vista previa).
\end{itemize}

\section{Ruta del capítulo}
Seis secciones con reglas claras, ejemplos mínimos y práctica corta.

\section{Uso de \textit{niet}}
\textbf{Para negar acciones y cualidades.} Usa \textit{niet} cuando niegas un verbo, un adjetivo, un adverbio, una preposición o un elemento específico de la frase.

\textbf{Ejemplos}
\begin{itemize}[leftmargin=*]
  \item \textit{Ik werk vandaag niet.} (Hoy no trabajo.)
  \item \textit{Het is niet duur.} (No es caro.)
  \item \textit{Hij woont niet in Utrecht.} (Él no vive en Utrecht.)
  \item \textit{We spreken niet snel.} (No hablamos rápido.)
\end{itemize}

\textbf{Mini práctica}
\begin{itemize}[leftmargin=*]
  \item Niega la acción: \textit{Ik kook vanavond}. $\rightarrow$ \textit{Ik kook vanavond niet.}
  \item Niega el lugar: \textit{Zij werkt in Amsterdam}. $\rightarrow$ \textit{Zij werkt niet in Amsterdam.}
\end{itemize}

\section{Uso de \textit{geen}}
\textbf{Para negar sustantivos indefinidos.} Emplea \textit{geen} cuando el sustantivo no lleva artículo definido o usa \textit{een}. Funciona con contables e incontables.

\textbf{Ejemplos}
\begin{itemize}[leftmargin=*]
  \item \textit{Ik heb geen auto.} (No tengo coche.)
  \item \textit{Wij drinken geen koffie.} (No bebemos café.)
  \item \textit{Zij is geen docent.} (Ella no es profesora.)
  \item Contraste: \textit{Ik heb niet de auto van Anna} (niega un coche concreto).
\end{itemize}

\textbf{Mini práctica}
\begin{itemize}[leftmargin=*]
  \item Transforma: \textit{Hij heeft een hond}. ➝ \textit{Hij heeft geen hond.}
  \item Niega un plural sin artículo: \textit{Wij hebben tickets}. ➝ \textit{Wij hebben geen tickets.}
\end{itemize}

\section{Posición de la negación en la frase}
\textbf{Regla A1:} coloca \textit{niet} casi al final, justo antes del elemento que quieres negar.

\textbf{Patrones útiles}
\begin{itemize}[leftmargin=*]
  \item Oración neutra: [Sujeto] + [V2] + \ldots{} + \textit{niet}.\ \textit{Ej.: Ik lees het boek niet}.
  \item Con complementos de lugar/tiempo: pon \textit{niet} antes de ese complemento.\ \textit{Ej.: Wij gaan niet naar de les vandaag}.
  \item Con infinitivo final: \textit{niet} va antes del infinitivo.\ \textit{Ej.: Ik wil niet werken}.
  \item Con participio final: \textit{Ik heb het huis niet gezien}. (Guía inicial; se refuerza en capítulos de pasado.)
\end{itemize}

\textbf{Mini práctica}
\begin{itemize}[leftmargin=*]
  \item Inserta \textit{niet}: \textit{Ik ga morgen naar Gent}. (Niega el viaje.)
  \item Mueve \textit{niet} para negar el lugar: \textit{Zij studeert in Leiden}.
\end{itemize}

\section{Negación con verbos}
\textbf{Verbo simple.} Coloca \textit{niet} al final del predicado verbal: \textit{Ik rook niet}.

\textbf{Modal + infinitivo.} \textit{niet} va antes del infinitivo final: \textit{Ik kan niet komen}. 

\textbf{Verbos separables.} Pon \textit{niet} antes de la partícula separada: \textit{Ik maak het raam niet open}. 

\textbf{Imperativo.} Forma breve: \textit{Niet roken!} / \textit{Niet doen!}

\textbf{Mini práctica}
\begin{itemize}[leftmargin=*]
  \item Niega con modal: \textit{Wij willen zwemmen}. ➝ \textit{Wij willen niet zwemmen}. 
  \item Con separable: \textit{Zij doet de lamp aan}. ➝ \textit{Zij doet de lamp niet aan}. 
\end{itemize}

\section{Negación con sustantivos}
\textbf{Indefinidos y cantidad cero.} Usa \textit{geen} para negar existencia o posesión cuando no hay artículo definido.

\begin{tabular}{@{}lll@{}}
\toprule
Afirmación & Negación & Traducción \\
\midrule
\textit{Hij heeft een fiets} & \textit{Hij heeft geen fiets} & Él no tiene bici \\
\textit{We hebben melk} & \textit{We hebben geen melk} & No tenemos leche \\
\textit{Er zijn stoelen} & \textit{Er zijn geen stoelen} & No hay sillas \\
\bottomrule
\end{tabular}

\textbf{Sustantivos definidos o posesivos.} No uses \textit{geen}. Niega con \textit{niet} porque hablas de algo concreto: \textit{Het is niet mijn boek}. 

\textbf{Mini práctica}
\begin{itemize}[leftmargin=*]
  \item Cambia a negación: \textit{Zij heeft een kat}. 
  \item Niega con posesivo: \textit{Dat is mijn fiets}. (Mantén \textit{niet}).
\end{itemize}

\section{Errores comunes}
\begin{itemize}[leftmargin=*]
\item Usar \textit{niet} con sustantivos indefinidos: *\textit{Ik heb niet een hond}. $\rightarrow$ \textit{Ik heb geen hond}.
  \item Colocar \textit{niet} demasiado pronto: *\textit{Ik niet ga}. Mantén el verbo en segunda posición: \textit{Ik ga niet}. 
  \item Olvidar la partícula separable: *\textit{Ik maak niet open}. Incluye el objeto: \textit{Ik maak de deur niet open}. 
  \item Doble negación innecesaria: *\textit{Ik heb geen niets}. Usa solo una: \textit{Ik heb niets}. 
  \item Mezclar con artículos definidos: \textit{niet} + artículo específico: \textit{Ik koop niet die jas}. 
\end{itemize}


\section{Práctica sugerida}
\begin{itemize}[leftmargin=*]
	\item Pares afirmación/negación transformando \textit{een} en \textit{geen}.
	\item Insertar \textit{niet} en la posición adecuada en frases dadas.
	\item Diálogos negando ofrecimientos o corrigiendo información.
\end{itemize}

\section{Microevaluación}
\begin{itemize}[leftmargin=*]
	\item 12 frases: 6 con \textit{geen}, 6 con \textit{niet}, revisadas en voz alta.
\end{itemize}
